% Junior Level 3, Buffer session B2 -- Free Build (Scratch or Python, child's choice)
\kcpart{B2}{Free Build}{Plan it, build it, test it, credit it: your project, your choice.}
\lessoninfo{Concepts: planning in parts (R), nested repeat and My Blocks (R), Python changes and \texttt{for} loops (R), testing and bug reports (R), remix and credit (R)}{Tools: Scratch 3 offline app, Thonny in simple mode (after week 6), idea cards, paper}{CSTA: 1B-AP-11, 1B-AP-12, 1B-AP-14, 1B-AP-15}

\begin{teachernote}[When to use B2]
Use B2 any time after week 5: as extra studio time before or after week 9, after a missed week, or as a reward. Before week 6, everyone builds in Scratch, because Python has not been met yet. From week 6 on, each child chooses Scratch or Python. If a Big Build is not finished, B2 is the perfect place to finish it (milestone M2) instead of starting something new.
\end{teachernote}

\begin{objectives}
  \item The student can choose Scratch or Python and plan a small project in two or three named parts before building. (J3.G2) [J3.B2.1]
  \item The student can build the project part by part, using a repeat inside a repeat or a My Block in Scratch, or changing numbers, colors and \texttt{for} loops in Python. (J3.G1, J3.G5) [J3.B2.2]
  \item The student can test a partner's project, describe a bug as \enquote{I expected \ldots, but it \ldots}, and fix a bug in their own project. (J3.G6) [J3.B2.3]
  \item The student can give credit in their project for an idea they remixed. (J3.G7) [J3.B2.4]
\end{objectives}

\begin{materials}
  \item The idea cards: six index cards you write by hand, one idea each: \enquote{a pattern that grows}, \enquote{a click-counting game}, \enquote{a turtle flag}, \enquote{a flower garden}, \enquote{a house with a sun}, \enquote{a mystery: my own idea}. Draw a little Scratch cat or a little turtle on each card to show which tool fits (the house and flag can be either).
  \item A half sheet of paper per child for a mini plan (three boxes: Part 1, Part 2, Part 3), and pencils.
  \item On each Mac: Scratch (with the Pen extension) and, after week 6, Thonny in simple mode, plus a \textbf{mouse}. Each child's earlier projects within reach: the week 1 and week 5 patterns, the week 8 Python programs, the Big Build.
  \item A kitchen timer (6--7 minutes) for Driver/Helper swaps. Optional: the Rose and Pixi video for this session.
\end{materials}

\needspace{16\baselineskip}
\section{At a glance}
\begin{timeplan}
  \timeslot{0--10}{Warm-up}{Idea cards: pick one or invent one. Draw a mini plan with two or three parts and choose Scratch or Python.}
  \timeslot{10--20}{Show me}{Remix and credit: start from an old project, save it with a new name, and add a credit line in Scratch and in Python.}
  \timeslot{20--33}{Make it}{Free build, part by part. Driver and Helper, swap on the timer.}
  \timeslot{33--35}{Wiggle break}{Freeze frame: act out your project in three poses.}
  \timeslot{35--50}{Make it}{Swap and test: partners try each other's projects and give one bug report. Fix, then keep building.}
  \timeslot{50--60}{Share}{Mini showcase with credits. Exit question and reflection.}
\end{timeplan}

\section{Warm-up (0--10 min)}
\begin{unplugged}[Idea cards and a mini plan]
\begin{enumerate}[leftmargin=1.4em,itemsep=1pt]
  \item Spread the idea cards on the table, face up. Each child picks one, or turns over the mystery card and says their own idea in one sentence. Give them 90 seconds.
  \item \enquote{Is it a Scratch idea or a Python idea?} Children circle the tool on their mini plan. A game that reacts to clicks or keys is a Scratch idea; a turtle picture can be either.
  \item Draw the mini plan: three boxes, one part each, in the order they will build them. Two parts is fine for a short session.
  \item If the idea comes from somewhere (a friend, an old project, a week of the course), write it under the boxes: \enquote{Idea from \ldots}
\end{enumerate}
With 1 child: pick an idea card too and plan your own project next to them; compare plans at the end. With 2--4: everyone tells the person next to them their parts in one sentence each.
\end{unplugged}

\section{Show me (10--20 min)}
\begin{teacherscript}
\enquote{Today is free build: you choose! You can start from nothing, or you can \textbf{remix}: start from an old project and change it into something new. Remixing is a great way to build. There is one rule: we \textbf{give credit}. We say whose idea we started from, right inside the project.}
\end{teacherscript}

\begin{analogy}[A recipe from Grandma]
\enquote{When you bake Grandma's cookies but add chocolate chips, you still say, \enquote{This is Grandma's recipe, with my chips.} A remix is the same: someone's idea, plus your change, plus a thank-you.}
\end{analogy}

\begin{workedexample}[Remix with a credit line]
Do both versions on your Mac while the children watch.
\begin{enumerate}[leftmargin=1.4em,itemsep=1pt]
  \item \textbf{Scratch.} Open the week 5 pattern. \menu{File} → \menu{Save to your computer} with a new name, such as \enquote{remix-garden}. Right-click an empty spot of the code area, choose \menu{Add Comment}, and type: \enquote{Remix of Kai's flower from Play Day.} Then change one thing: the pen color.
  \item \textbf{Python.} In Finder, copy a week 8 program (for example the two-color flower's \texttt{starter.py}) into \texttt{my-python} and rename the copy with a clear name, such as \texttt{remix\_flower.py}. In Thonny, click \menu{Load} and open the copy. Find the \texttt{\# Changed by:} line at the top and type a nickname after it, or add a new line \texttt{\# Remix of the course's two-color flower}. \enquote{A line that starts with \texttt{\#} is a note for people. Python skips it.} Change one color, then press Run.
\end{enumerate}
\enquote{Saving with a new name keeps the old project safe. The credit line says thank you.}
\end{workedexample}

\section{Make it (20--33 min)}
\begin{pairroles}[Driver and Helper on two projects]
Each child builds their own project. The \textbf{Driver} builds. The \textbf{Helper} holds the Driver's mini plan, reads the next part, and says \enquote{Test it!} after every change; the Helper never grabs the mouse or the keyboard. Swap on the timer (6--7 minutes) so each child drives their own project. With 1 child: you are the Helper the whole time. With 3: two Drivers and one Helper who moves between them on the timer. With 4: two pairs.
\end{pairroles}

\begin{freebuild}[Your project, your choice]
Children build their mini plan, one part at a time, testing after each part. Walk around and ask: \enquote{Which part are you on? Show me it working.} Ideas for children who finish early: add a sound or a counter (Scratch); add a second color or a thicker pen (Python); turn a repeated shape into a My Block or a \texttt{def} (stretch). Remind Python builders: change one thing, then Run.
\end{freebuild}

\section{Wiggle break (33--35 min)}
Freeze frame: each child acts out their project in three poses, one for each part, while the others guess what it is. Then shake out arms and legs.

\section{Make it (35--50 min)}
\begin{bughunt}[Swap and test (about 6 minutes)]
At about minute 38, partners swap seats. The visitor holds the owner's mini plan, runs the project and tries everything. Then they give \textbf{one} bug report, in the shape from week 7: \enquote{I expected \ldots, but it \ldots}, for example \enquote{I expected the flower to be gold, but it was purple} or \enquote{I expected the counter to start at 0, but it started at 5.} If they find no bug, they give one \enquote{I liked \ldots}. Swap back: the owner fixes their own bug, testing after the fix. Remind them: the visitor never changes someone else's project.
\end{bughunt}

After the fix, children keep building until minute 50. Save often.

\section{Share (50--60 min)}
\begin{safety}[Credit and nicknames]
\enquote{Show us where your credit line is.} Each child points to their Scratch comment or their Python \texttt{\#} line. Then talk about two rules:
\begin{itemize}[leftmargin=1.2em,itemsep=0pt]
  \item \textbf{Give credit.} Copying a friend's idea and saying it is all yours is not fair. Remixing and saying thank you is.
  \item \textbf{Nicknames in credit lines.} \enquote{Remix of Kai's flower} is fine. Never put a full name, school or address in a project, even in a comment.
\end{itemize}
\end{safety}

\begin{teacherscript}
Mini showcase: each child runs their project and says one sentence about one part of their code. Others give a thumbs up and one \enquote{I liked \ldots}.
\end{teacherscript}

\begin{formativecheck}[Exit question]
\textbf{Look for:} a mini plan with two or three parts; each part tested; one bug report given and one bug fixed; a credit line where an idea was remixed.

\textbf{Ask:} \enquote{Why do we save a remix with a new name?}
\textbf{Answer:} so the old project stays safe and does not get changed.
\end{formativecheck}

\begin{reflection}
Ask: \enquote{Which part of today's project are you proudest of?} and \enquote{Did you choose Scratch or Python, and why?} Children draw a smiley on their mini plan next to the part they like best.
\end{reflection}

\section{Notes for the teacher}
\begin{commonmistakes}
  \item \textbf{A plan too big for one session.} Two parts that work beat five parts that do not. Extra ideas go on the back of the plan.
  \item \textbf{Saving over the old project.} In Scratch, use \menu{Save to your computer} with a new name before changing anything. For Python, copy the file in Finder first (Thonny's simple mode has no \enquote{Save as}), then load the copy.
  \item \textbf{A Python comment that breaks the program.} A note without \texttt{\#} at the start gives a red message. Put \texttt{\#} at the very start of the line.
  \item \textbf{Bug reports that are only \enquote{it's broken}.} Model the sentence: \enquote{I expected \ldots, but it \ldots}
  \item \textbf{Stuck on choosing.} Give two choices only: \enquote{the flower garden in Scratch, or the house in Python?}
\end{commonmistakes}

\begin{differentiation}
\supportpath Start from a finished project and change only numbers and colors: the week 5 pattern in Scratch or a week 8 program in Python. A one-part plan is fine. Write the credit line together.
\extensionpath Build the same picture in both tools: a Scratch pattern and its Python twin, using the blocks-to-words cards from week 5. Or use a My Block or a \texttt{def} for a shape you draw three times.
\end{differentiation}

\begin{stretchchallenge}[Twin projects]
Pick one small part of your project and build it in the other tool too. Does the Python square look the same as the Scratch square? Which numbers had to match?
\end{stretchchallenge}

\begin{teachernote}[If you are short on time]
Skip the idea cards and let children continue their Big Build or a favorite earlier project. Keep the swap and test, even if short: it is this session's bug hunt. Keep the credit talk at the Share.
\end{teachernote}
